%%
% BIThesis 本科毕业设计论文模板 —— 使用 XeLaTeX 编译 The BIThesis Template for Undergraduate Thesis
% This file has no copyright assigned and is placed in the Public Domain.
%%

\begin{conclusion}
本文围绕“基于微服务架构的在线视频分享平台的设计与实现”展开研究与开发，完成了一个具备视频浏览、视频搜索、在线播放、用户互动、视频投稿、分片上传、视频转码、后台审核和 AI 辅助审核等功能的在线视频分享平台原型。系统以视频内容生命周期为主线，围绕视频从用户投稿到正式发布的完整过程，设计并实现了上传、处理、审核、发布和播放等关键环节，基本满足在线视频分享平台的核心业务需求。

在系统架构方面，本文采用微服务架构对后端功能进行拆分，将网关服务、主站业务服务、资源文件服务、互动服务、后台管理服务和公共基础模块分别承担不同职责，降低了系统内部模块耦合度。系统通过 Nacos 实现服务注册与发现，通过 Spring Cloud Gateway 提供统一访问入口，通过 OpenFeign 实现服务间调用，并结合 Redis、Elasticsearch、RabbitMQ 和 Seata 等组件完成缓存、检索、异步任务和跨服务一致性控制。该架构设计使系统具有较清晰的服务边界和较好的扩展基础。

在视频处理方面，本文实现了大文件分片上传、文件合并、FFmpeg 视频转码和 HLS 播放资源生成。该流程能够降低大文件上传失败后的重传成本，并将用户上传的原始视频转换为适合网页端播放的流媒体资源。在内容审核方面，本文设计了“AI 初审 + 人工终审”的协同审核机制。AI 审核服务通过消息队列接收审核任务，对视频语音、声音事件和关键帧内容进行分析，并返回风险等级、风险标签和审核建议。管理员在后台结合 AI 审核结果进行最终判断，从而在提高审核效率的同时保留人工决策能力。

在系统测试方面，本文对用户端功能、创作中心功能、后台管理功能、视频处理流程、AI 审核流程、跨服务事务一致性和异常场景进行了测试。测试结果表明，系统能够完成视频投稿、上传、转码、AI 审核、人工审核、正式发布和前台播放等核心流程，用户互动、搜索检索、后台管理和异常处理等功能也能够正常运行，整体业务链路能够形成闭环。

虽然本文完成了系统核心功能的设计与实现，但系统仍存在一些不足。首先，系统主要在本地和小规模环境下完成测试，尚未进行大规模并发压力测试；其次，视频文件主要采用本地目录存储，在多实例部署和大规模文件管理场景下仍存在扩展性不足的问题；再次，AI 审核结果受模型能力、视频质量和上下文语义影响，仍可能出现误判或漏判，因此目前更适合作为人工审核的辅助工具；最后，系统自动化测试、日志监控和持续集成能力仍有进一步完善空间。

后续工作可以从以下几个方面继续优化：一是引入对象存储或分布式文件系统，提高视频资源存储和访问能力；二是对首页访问、视频播放、评论弹幕、搜索查询和视频上传等高频场景进行压力测试，并根据测试结果优化缓存策略和服务扩容方案；三是进一步完善 AI 审核规则体系和风险标签体系，结合人工反馈提升审核准确率和可解释性；四是补充接口自动化测试、单元测试、日志追踪和系统监控，提高系统长期运行和维护能力。

综上所述，本文完成了一个基于微服务架构的在线视频分享平台原型，实现了视频处理、内容审核、用户互动和后台管理等核心功能。该系统能够体现微服务架构、异步任务处理、视频转码播放和智能内容审核技术在复杂 Web 平台中的综合应用价值，也为后续进一步扩展为更完善的视频内容服务平台奠定了基础。
\end{conclusion}